\documentclass{article}
\usepackage{comment}
\usepackage{amsmath}
\usepackage{amsfonts}
\usepackage{sectsty}
\usepackage{hyperref}
\usepackage{fancyhdr}
\usepackage[top=1in,bottom=1in,left=1in,right=1in]{geometry}
\usepackage{float}
\usepackage{graphicx}
\usepackage{tabularx}
\usepackage{mathtools}

\title{Lecture 3 - The Impulse Response}
\author{Matthew Farah, \href{mailto:farahm11@mcmaster.ca}{$\langle$farahm11@mcmaster.ca$\rangle$}, 400468251}
\date{\today}

\hypersetup{
        colorlinks=true,
        linkcolor=blue,
        filecolor=magenta,
        urlcolor=blue,
        citecolor=blue,
	anchorcolor=blue
}

\fancyhead{}
\fancyhead[L]{Matthew Farah, \href{mailto:farahm11@mcmaster.ca}{$\langle$farahm11@mcmaster.ca$\rangle$}, 400468251}
\fancyhead[R]{Lecture 3 - The Impulse Response}

\newenvironment{directsubsubs}{
  \makeatletter
  \renewcommand\thesubsubsection{\thesection.\Roman{subsubsection}}
  \makeatother
}{
  \makeatletter
  \renewcommand\thesubsubsection{\thesubsection.\Roman{subsubsection}}
  \makeatother
}

\renewcommand{\thesubsection}{\thesection.\alph{subsection}}
\renewcommand{\thesubsubsection}{\thesection.\alph{subsection}.\Roman{subsubsection}}
\makeatletter
\DeclareFontFamily{U}{MnSymbolA}{}
\DeclareFontShape{U}{MnSymbolA}{m}{n}{
    <-6>  MnSymbolA5
   <6-7>  MnSymbolA6
   <7-8>  MnSymbolA7
   <8-9>  MnSymbolA8
   <9-10> MnSymbolA9
  <10-12> MnSymbolA10
  <12->   MnSymbolA12}{}
\DeclareFontShape{U}{MnSymbolA}{b}{n}{
    <-6>  MnSymbolA-Bold5
   <6-7>  MnSymbolA-Bold6
   <7-8>  MnSymbolA-Bold7
<8-9>  MnSymbolA-Bold8
   <9-10> MnSymbolA-Bold9
  <10-12> MnSymbolA-Bold10
  <12->   MnSymbolA-Bold12}{}
\DeclareSymbolFont{MnSyA}{U}{MnSymbolA}{m}{n}
\SetSymbolFont{MnSyA}{bold}{U}{MnSymbolA}{b}{n}

\DeclareRobustCommand{\overleftharpoon}{\mathpalette{\overarrow@\leftharpoonfill@}}
\DeclareRobustCommand{\overrightharpoon}{\mathpalette{\overarrow@\rightharpoonfill@}}
\def\leftharpoonfill@{\arrowfill@\leftharpoondown\mn@relbar\mn@relbar}
\def\rightharpoonfill@{\arrowfill@\mn@relbar\mn@relbar\rightharpoonup}

\DeclareMathSymbol{\leftharpoondown}{\mathrel}{MnSyA}{'112}
\DeclareMathSymbol{\rightharpoonup}{\mathrel}{MnSyA}{'100}
\DeclareMathSymbol{\mn@relbar}{\mathrel}{MnSyA}{'320}
\makeatother


\newcolumntype{Y}{>{\centering\arraybackslash}X}
\renewcommand{\arraystretch}{1.8}

\sectionfont{\fontsize{12}{12}\bfseries}
\subsectionfont{\fontsize{10}{12}\normalfont}
\subsubsectionfont{\fontsize{10}{12}\normalfont}

\newcommand{\harpoon}[1]{\overrightharpoonup{#1}}

\begin{document}
\maketitle
\pagestyle{fancy}

\begin{itemize}
	\item Systems can be represented:
	\begin{enumerate}
		\item As Difference Equations (DEQs).
		\item In State-Space Form.
		\begin{itemize}
			\item The state-space form of a system is analogous to an infinite state machine
		\end{itemize}
	\end{enumerate}
\end{itemize}

\begin{directsubsubs}
\section[Example ~\thesection]{Represent the system $y(n) = x(n) + x(n-1) - y(n-1)$}

\subsubsection[~\thesubsubsection]{Find $\overrightharpoon{S(n)}$}

\begin{align*}
	&\text{Let:} && \text{Then:} \\
	&s_1(n) = x(n-1) && s_1(n+1) = x(n) \\
	&s_2(n) = y(n-1) && s_2(n+1) = y(n) = x(n) + x(n-1) - y(n-1) = x(n) + s_1(n) - s_2(n)
\end{align*}

Thus, let $\overrightharpoon{S(n)} = \begin{psmallmatrix} s_1(n) \\ s_2(n) \end{psmallmatrix}$.

\subsubsection[~\thesubsubsection]{Find $\overrightharpoon{S(n+1)}$}

\begin{align*}
	\overrightharpoon{S(n+1)} &= \begin{pmatrix} s_1(n+1) \\ s_2(n+1) \end{pmatrix} \\
	&= \begin{pmatrix} x(n) \\ x(n) + s_1(n) - s_2(n) \end{pmatrix} \\
	&= \begin{pmatrix} 0 & 0 \\ 1 & -1 \end{pmatrix}\overrightharpoon{S(n)} + \begin{pmatrix} 1 \\ 1 \end{pmatrix}x(n)
\end{align*}

Therefore, our $\boldsymbol{A}$ and $\overrightharpoon{B}$ matrices are $\begin{psmallmatrix} 0 & 0 \\ 1 & -1 \end{psmallmatrix}$ and $\begin{psmallmatrix} 1 \\ 1 \end{psmallmatrix}$ respectively.

\subsubsection[~\thesubsubsection]{Find $y(n)$}

\begin{align*}
	y(n) &= x(n) + x(n-1) - y(n-1) \\
	&= s_1(n) - s_2(n) + x(n) \\
	&= \begin{pmatrix} -1 & 1 \end{pmatrix}\overrightharpoon{S(n)} + \begin{pmatrix} 1 \end{pmatrix}x(n)
\end{align*}

Therefore, our $\boldsymbol{C}$ and $\boldsymbol{D}$ matrices are $\begin{psmallmatrix} -1 & 1 \end{psmallmatrix}$ and $\begin{psmallmatrix} 1 \end{psmallmatrix}$ respectively.

Putting it all together, our system can be represented in state space form as:

\begin{align*}
	\overrightharpoon{S(n+1)} &= \begin{pmatrix} 0 & 0 \\ 1 & -1 \end{pmatrix}\overrightharpoon{S(n)} + \begin{pmatrix} 1 \\ 1 \end{pmatrix}x(n) \\ \\
	y(n) &= \begin{pmatrix} 0 & 0 \\ 1 & -1 \end{pmatrix}\overrightharpoon{S(n)} + \begin{pmatrix} 1 \\ 1 \end{pmatrix}x(n)
\end{align*}

\begin{itemize}
	\item The Kronecker-Delta function $\delta(n - k)$ is defined as:
	\begin{itemize}
		\item $\delta(n - k) = \begin{cases} 1, \; \text{if} \; n = k \\ 0 \; \text{otherwise} \end{cases}$
		\item The Kronecker-Delta function acts as the discrete analog to the continuous Dirac-Delta function.
	\end{itemize}
\end{itemize}

\end{directsubsubs}

\section[Example ~\thesection]{Using the state space equations from example 1, compute the outputs to the system for the input $x(n) \delta(n - 2)$}

\begin{align*}
	&\underline{\text{Current State} \; (\overrightharpoon{S(n)})}& &\underline{\text{Current Output} \; (y(n))} \\ \\
	&\overrightharpoon{S(0)} = \overrightharpoon{0}& &y(0) = \begin{pmatrix} 1 & -1 \end{pmatrix}\overrightharpoon{S(0)} + \begin{pmatrix} 1 \end{pmatrix}x(0) = 0 \\
	&\overrightharpoon{S(1)} = \begin{pmatrix} 0 & 0 \\ 1 & -1 \end{pmatrix}\overrightharpoon{S(0)} + \begin{pmatrix} 1 \\ 1 \end{pmatrix}x(0) = \overrightharpoon{0}& &y(1) = \begin{pmatrix} 1 & -1 \end{pmatrix}\overrightharpoon{S(1)} + \begin{pmatrix} 1 \end{pmatrix}x(1) = 0 \\
	&\overrightharpoon{S(2)} = \begin{pmatrix} 0 & 0 \\ 1 & -1 \end{pmatrix}\overrightharpoon{S(1)} + \begin{pmatrix} 1 \\ 1 \end{pmatrix}x(1) = \overrightharpoon{0}& &y(2) = \begin{pmatrix} 1 & -1 \end{pmatrix}\overrightharpoon{S(2)} + \begin{pmatrix} 1 \end{pmatrix}x(2) = 1 \\
	&\overrightharpoon{S(3)} = \begin{pmatrix} 0 & 0 \\ 1 & -1 \end{pmatrix}\overrightharpoon{S(2)} + \begin{pmatrix} 1 \\ 1 \end{pmatrix}x(2) = \begin{pmatrix} 1 \\ 1 \end{pmatrix}& &y(3) = \begin{pmatrix} 1 & -1 \end{pmatrix}\overrightharpoon{S(3)} + \begin{pmatrix} 1 \end{pmatrix}x(3) = 0 \\
	&\overrightharpoon{S(4)} = \begin{pmatrix} 0 & 0 \\ 1 & -1 \end{pmatrix}\overrightharpoon{S(3)} + \begin{pmatrix} 1 \\ 1 \end{pmatrix}x(3) = \overrightharpoon{0} & &y(4) = \begin{pmatrix} 1 & -1 \end{pmatrix}\overrightharpoon{S(4)} + \begin{pmatrix} 1 \end{pmatrix}x(4) = 0 \\
	&\vdots & &\vdots
\end{align*}

\section[Example ~\thesection]{Model the system $y(n) = x(n) + {\alpha}y(n - 1)$ in state-space form and compute its outputs for $x(n) = \delta(n)$.}
Note that this system models compound interest.

\subsection[~\thesubsection]{Express system in state-space form}

\subsubsection[\thesubsubsection]{Find $S(n)$}

\begin{align*}
	&\text{Let:}& &\text{Then:} \\
	&s_1(n) - y(n - 1)& &s_1(n + 1) = y(n) = x(n) + {\alpha}y(n - 1) = {\alpha}s_1(n) + x(n) \\
\end{align*}

Thus, let $\overrightharpoon{S(n)} = \begin{psmallmatrix} s_1(n) \end{psmallmatrix}$

\subsubsection[~\thesubsubsection]{Find $S(n+1)$}

\begin{align*}
	\overrightharpoon{S(n+1)} &= \begin{pmatrix} s_1(n + 1) \end{pmatrix} \\
	&= \begin{pmatrix} {\alpha}s_1(n) + x(n) \end{pmatrix} \\
	&= \begin{pmatrix} \alpha \end{pmatrix}\overrightharpoon{S(n)} + \begin{pmatrix} 1 \end{pmatrix}x(n)
\end{align*}

Therefore, our $\boldsymbol{A}$ and $\overrightharpoon{B}$ matrices are $\begin{psmallmatrix} \alpha \end{psmallmatrix}$ and $\begin{psmallmatrix} 1 \end{psmallmatrix}$ respectively.

\subsubsection[~\thesubsubsection]{Find $y(n)$}

\begin{align*}
	y(n) &= x(n) + {\alpha}y(n - 1) \\
	&= s_1(n) + x(n) \\
	&= \begin{pmatrix} \alpha \end{pmatrix}\overrightharpoon{S(n)} + \begin{pmatrix} 1 \end{pmatrix}x(n)
\end{align*}

Therefore, our $\boldsymbol{C}$ and $\boldsymbol{D}$ matrices are $\begin{psmallmatrix} \alpha \end{psmallmatrix}$ and $\begin{psmallmatrix} 1 \end{psmallmatrix}$ respectively.

Putting it all together, our system can be represented in state space form as:

\begin{align*}
	\overrightharpoon{S(n + 1)}&= \begin{pmatrix} \alpha \end{pmatrix}\overrightharpoon{S(n)} + \begin{pmatrix} 1 \end{pmatrix}x(n) \\
	y(n) &= \begin{pmatrix} \alpha \end{pmatrix}\overrightharpoon{S(n)} + \begin{pmatrix} 1 \end{pmatrix}x(n)
\end{align*}

\subsection[~\thesubsection]{Compute the inputs of $y(n)$}

\begin{align*}
	&\underline{\text{Current State} \; (\overrightharpoon{S(n)})}& &\underline{\text{Current Output} \; (y(n))} \\ \\
	&\overrightharpoon{S(0)} = \overrightharpoon{0}& &y(0) = \begin{pmatrix} \alpha \end{pmatrix}\overrightharpoon{S(0)} + \begin{pmatrix} 1 \end{pmatrix}x(0) = 1 \\
	&\overrightharpoon{S(1)} = \begin{pmatrix} \alpha \end{pmatrix}\overrightharpoon{S(0)} + \begin{pmatrix} 1 \end{pmatrix}x(0) = 1 & &y(1) = \begin{pmatrix} \alpha \end{pmatrix}\overrightharpoon{S(1)} + \begin{pmatrix} 1 \end{pmatrix}x(1) = \alpha \\
	&\overrightharpoon{S(2)} = \begin{pmatrix} \alpha \end{pmatrix}\overrightharpoon{S(1)} + \begin{pmatrix} 1 \end{pmatrix}x(1) = \alpha & &y(2) = \begin{pmatrix} \alpha \end{pmatrix}\overrightharpoon{S(2)} + \begin{pmatrix} 1 \end{pmatrix}x(2) = {\alpha}^2 \\
	&\overrightharpoon{S(3)} = \begin{pmatrix} \alpha \end{pmatrix}\overrightharpoon{S(2)} + \begin{pmatrix} 1 \end{pmatrix}x(2) = {\alpha}^2 & &y(3) = \begin{pmatrix} \alpha \end{pmatrix}\overrightharpoon{S(3)} + \begin{pmatrix} 1 \end{pmatrix}x(3) = {\alpha}^3 \\
	&\vdots& &\vdots
\end{align*}

Generalizing, for input $x(n) = \delta(n)$, $y(n) = {\alpha}^n$ for this system.

\section[Example ~\thesection]{Describe the outputs of any system for the input $x(n) = \delta(n - 1)$ using the state-space form of systems}

We know that any system can be represented in state space form by the equations:

\begin{align*}
	\overrightharpoon{S(n+1)} &= \boldsymbol{A}\overrightharpoon{S(n)} + \overrightharpoon{B(n)}x(n) \\
	y(n) &= \boldsymbol{C}\overrightharpoon{S(n)} + \boldsymbol{D}x(n)
\end{align*}

For the input $x(n) = \delta(n)$, we can deduce that:

\begin{align*}
	&\underline{\text{Current State} \; (\overrightharpoon{S(n)})}& &\underline{\text{Current Output} \; (y(n))} \\ \\
	&\overrightharpoon{S(0)} = \overrightharpoon{0}& &y(0) = \boldsymbol{C}\overrightharpoon{S(0)} + \boldsymbol{D}x(0) = \boldsymbol{D} \\
	&\overrightharpoon{S(1)} = \boldsymbol{A}\overrightharpoon{S(0)} + \overrightharpoon{B}x(0) = \overrightharpoon{B}& & y(1) = \boldsymbol{C}\overrightharpoon{S(1)} + \boldsymbol{D}x(1) = \boldsymbol{C}\overrightharpoon{B} \\
	&\overrightharpoon{S(2)} = \boldsymbol{A}\overrightharpoon{S(1)} + \overrightharpoon{B}x(1) = \boldsymbol{A}\overrightharpoon{B} & &y(2) = \boldsymbol{C}\overrightharpoon{S(2)} + \boldsymbol{D}x(2) = \boldsymbol{C}\boldsymbol{A}\overrightharpoon{B} \\
	&\overrightharpoon{S(3)} = \boldsymbol{A}\overrightharpoon{S(2)} + \overrightharpoon{B}x(2) = \boldsymbol{A}^2\overrightharpoon{B} & &y(3) = \boldsymbol{C}\overrightharpoon{S(3)} + \boldsymbol{D}x(3) = \boldsymbol{C}\boldsymbol{A}^2\overrightharpoon{B} \\
	&\vdots& &\vdots
\end{align*}

Thus, we can describe the inputs to any system with input $x(n) = \delta(n)$ in state-space form as:

\begin{align*}
	y(n) &= \begin{cases} 0, \; \text{if} \; n < 0 \\ \boldsymbol{D}, \text{if} \; n = 0 \\ \boldsymbol{C}\boldsymbol{A}^{n-1}\overrightharpoon{B} \end{cases}
\end{align*}

\begin{itemize}
	\item The impulse response of a system is defined as the outputs of a system when $x(n) = \delta(n)$
	\begin{itemize}
		\item Denoted $h(n)$
	\end{itemize}
\end{itemize}

\section[Example ~\thesection]{Represent the system $y(n) = {\alpha_1}x(n) + {\alpha_2}x(n - 1) + {\alpha_3}x(n - 2)$ in state-space form and describe its impulse response}

Note that we can clearly take a shortcut here and realize that for $x(n) = \delta(n)$, $y(n) = {\alpha_1}\delta(n) + {\alpha_2}\delta(n - 1) + {\alpha_3}\delta(n - 2)$, and thus:

\begin{align*}
	h(n) &= \begin{cases} \alpha_1, \; \text{if} \; n = 0 \\ \alpha_2, \; \text{if} n = 1 \\ \alpha_3, \; \text{if} \; n = 2 \\ 0, \; \text{otherwise} \end{cases}
\end{align*}

However, we will do as the question asks

\subsection[~\thesubsection]{Represent the system in state-space form}

\subsubsection[~\thesubsubsection]{Find $\overrightharpoon{S(n)}$}

\begin{align*}
	&\text{Let:}& &\text{Then:} \\
	&s_1(n) = x(n - 2) & &s_1(n + 1) = x(n - 1) = s_2(n) \\
	&s_2(n) = x(n - 1) & &s_2(n + 1) = x(n)
\end{align*}

Thus, let $\overrightharpoon{S(n)} = \begin{psmallmatrix} s_1(n) \\ s_2(n) \end{psmallmatrix}$.

\subsubsection[~\thesubsubsection]{Find $S(n + 1)$}

\begin{align*}
	\overrightharpoon{S(n + 1)} &= \begin{pmatrix} s_1(n + 1) \\ s_2(n + 1) \end{pmatrix} \\
	&= \begin{pmatrix} s_2(n) \\ x(n) \end{pmatrix} \\
	&= \begin{pmatrix} 0 & 1 \\ 0 & 0 \end{pmatrix}\overrightharpoon{S(n)} + \begin{pmatrix} 0 \\ 1 \end{pmatrix}x(n)
\end{align*}

Therefore, our $\boldsymbol{A}$ and $\overrightharpoon{B}$ are $\begin{psmallmatrix} 0 \\ 1 \end{psmallmatrix}$ and $\begin{psmallmatrix} 0 \\ 1 \end{psmallmatrix}$ respectively.

\subsubsection[~\thesubsubsection]{Find $y(n)$}

\begin{align*}
	y(n) &= {\alpha_1}x(n) + {\alpha_2}x(n - 1) + {\alpha_3}x(n - 2) \\
	&= {\alpha_3}s_1(n) + {\alpha_2}s_2(n) + {\alpha_1}x(n) \\
	&= \begin{pmatrix} \alpha_3 & \alpha_2 \end{pmatrix}\overrightharpoon{S(n)} + \begin{pmatrix} \alpha_1 \end{pmatrix}x(n)
\end{align*}

Therefore, our $\boldsymbol{C}$ and $\boldsymbol{D}$ matrices are $\begin{psmallmatrix} \alpha_3 & \alpha_2 \end{psmallmatrix}$ and $\begin{psmallmatrix} \alpha_1 \end{psmallmatrix}$ respectively.

Putting it all together, our system can be represented in state space form as:

\begin{align*}
	\overrightharpoon{S(n + 1)} &= \begin{pmatrix} 0 & 1 \\ 0 & 0 \end{pmatrix}\overrightharpoon{S(n)} + \begin{pmatrix} 0 \\ 1 \end{pmatrix}x(n) \\
	y(n) &= \begin{pmatrix} \alpha_3 & \alpha_2 \end{pmatrix}\overrightharpoon{S(n)} + \begin{pmatrix} \alpha_1 \end{pmatrix}x(n)
\end{align*}

\subsection[~\thesubsection]{Compute $y(n)$ for $x(n) = \delta(n)$ using the state space equations}

\begin{align*}
	&\underline{\text{Current State} \; (\overrightharpoon{S(n)})}& &\underline{\text{Current Output} \; (y(n))} \\ \\
	&\overrightharpoon{S(0)} = \overrightharpoon{0} & &y(0) = \begin{pmatrix} \alpha_3 & \alpha_2 \end{pmatrix}\overrightharpoon{S(0)} + \begin{pmatrix} \alpha_1 \end{pmatrix}x(0) = \alpha_1 \\
	&\overrightharpoon{S(1)} = \begin{pmatrix} 0 & 1 \\ 0 & 0 \end{pmatrix}\overrightharpoon{S(0)} + \begin{pmatrix} 0 \\ 1 \end{pmatrix}x(0) = \begin{pmatrix} 0 \\ 1 \end{pmatrix}& &y(1) = \begin{pmatrix} \alpha_3 & \alpha_2 \end{pmatrix}\overrightharpoon{S(1)} + \begin{pmatrix} \alpha_1 \end{pmatrix}x(1) = \alpha_2 \\
	&\overrightharpoon{S(2)} = \begin{pmatrix} 0 & 1 \\ 0 & 0 \end{pmatrix}\overrightharpoon{S(1)} + \begin{pmatrix} 0 \\ 1 \end{pmatrix}x(1) = \begin{pmatrix} 1 \\ 0 \end{pmatrix}& &y(2) = \begin{pmatrix} \alpha_3 & \alpha_2 \end{pmatrix}\overrightharpoon{S(2)} + \begin{pmatrix} \alpha_1 \end{pmatrix}x(2) = \alpha_3 \\
	&\overrightharpoon{S(3)} = \begin{pmatrix} 0 & 1 \\ 0 & 0 \end{pmatrix}\overrightharpoon{S(2)} + \begin{pmatrix} 0 \\ 1 \end{pmatrix}x(2) = \overrightharpoon{0}& &y(3) = \begin{pmatrix} \alpha_3 & \alpha_2 \end{pmatrix}\overrightharpoon{S(3)} + \begin{pmatrix} \alpha_1 \end{pmatrix}x(3) = 0 \\
	&\vdots& &\vdots
\end{align*}

Thus, the impulse response of the system is given by:

\begin{align*}
	h(n) &= \begin{cases} \alpha_1 \; \text{if} \; n = 0 \\ \alpha_2 \; \text{if} \; n = 1 \\ \alpha_3 \; \text{if} \; n = 2 \\ 0 \; \text{otherwise} \end{cases}
\end{align*}

\begin{itemize}
	\item A \emph{Finite Impulse Response} (FIR) system is a system for where there exists some same after which the system's output is always equal to zero.
		\begin{itemize}
			\item Expressed in mathematical notation, for an FIR system, $\exists N \in \mathbb{N} \quad \text{such that} \quad \forall n \ge N, h(n) = 0$
			\item FIR systems have a finite number of possible outputs for its impulse response.
			\item FIR systems can be expressed as a linear combination of Kronecker-Delta functions.
		\end{itemize}
	\item An \emph{Infinite Impulse Response} (IIR) system is any system which is \underline{NOT} an FIR system.
	\item All systems without feedback are FIR systems, but not all systems with feedback are IIR.
\end{itemize}

\end{document}
